\part{Business plan}
\chapter{Introduzione al Business Plan}
Il Business Plan è uno strumento fondamentale per chiunque voglia avviare o sviluppare un’attività, che si tratti di una startup, di un’azienda consolidata o di un progetto innovativo. Sebbene non sia obbligatorio redigerlo, è altamente consigliato poiché rappresenta una guida per mantenere il focus e per definire i passi necessari verso il successo.

\begin{redbar}
Il \textbf{Business Plan}\index{Business plan} permette di stabilire obiettivi chiari, pianificare strategie efficaci e valutare la fattibilità economica di un progetto imprenditoriale. Questo documento può essere immaginato come una mappa: delinea il percorso da seguire, segnalando le tappe fondamentali e aiutando a prevedere gli ostacoli che potrebbero emergere lungo il cammino.
\end{redbar}

\section{A chi è rivolto il Business Plan?}
Il Business Plan ha molteplici destinatari, ciascuno con esigenze e obiettivi specifici. 
\begin{itemize}
    \item Per gli \textit{imprenditori e i fondatori} di startup, è uno strumento per chiarire l’idea imprenditoriale, pianificare strategie di crescita e monitorare i progressi. Inoltre, consente di identificare i rischi e adattarsi ai cambiamenti che inevitabilmente si presenteranno durante il percorso.
    \item Per \textit{investitori e banche}, invece, il Business Plan è cruciale per valutare la fattibilità economica di un progetto, stimare i ritorni sull’investimento e comprendere i rischi associati. In particolare, gli istituti di credito utilizzano questo documento per decidere se concedere prestiti o linee di credito.
    \item Anche i \textit{partner commerciali} beneficiano del Business Plan, poiché fornisce loro una panoramica della solidità del progetto e degli obiettivi aziendali, elementi essenziali prima di avviare collaborazioni. 
    \item Allo stesso modo, per i \textit{membri del team e i dipendenti}, il Business Plan è una guida che aiuta a comprendere la visione e gli obiettivi aziendali, favorendo l’allineamento con le strategie aziendali.
    \item Non meno importanti sono i \textit{consulenti e i mentor}, che utilizzano il Business Plan per orientare gli imprenditori nelle decisioni strategiche e offrire suggerimenti per migliorare il piano. 
    \item Infine, \textit{enti pubblici} e organizzazioni di finanziamento lo utilizzano per valutare la possibilità di concedere fondi, agevolazioni fiscali o altri incentivi, nonché per garantire che il progetto rispetti i requisiti normativi.
\end{itemize}

\section{Struttura del Business Plan}
Un Business Plan si compone di due principali aspetti: uno qualitativo e uno quantitativo. La descrizione qualitativa precede sempre quella quantitativa, costituendo il presupposto fondamentale per la costruzione di un bilancio previsionale. La componente \textit{qualitativa} permette di descrivere le iniziative che l’impresa intende realizzare, mentre quella \textit{quantitativa} traduce queste idee in numeri e proiezioni economiche.

L'obiettivo di questa struttura è ridurre la soggettività nel processo decisionale e migliorare la comunicazione esterna con stakeholder, partner e potenziali investitori. In questo modo, il Business Plan diventa un ponte tra l’idea imprenditoriale e la sua realizzazione concreta.

Un Business Plan completo si articola in diverse sezioni, ciascuna delle quali svolge un ruolo specifico. 
\begin{enumerate}
    \item La prima è l’\textit{Executive Summary}, una sintesi del progetto che cattura l’essenza del piano e ne evidenzia i punti chiave. Questo documento deve essere chiaro, conciso e convincente, in modo da attrarre immediatamente l’attenzione dei lettori.
    \item Segue l’\textit{analisi del mercato di riferimento}, una sezione in cui si descrive il contesto in cui l’impresa opererà, individuando i segmenti di clientela, le tendenze e le opportunità. 
    \item Subito dopo, si passa a identificare il \textit{problema} che il prodotto o servizio intende risolvere, ponendo le basi per la successiva sezione dedicata alla soluzione proposta.
    \item La \textit{Unique Selling Proposition (USP)} è un elemento chiave, in quanto definisce cosa rende l’offerta unica rispetto alla concorrenza. 
    \item A tal proposito, l’\textit{analisi dei competitor e dei vantaggi competitivi}, insieme alla \textit{SWOT analysis}, permette di valutare punti di forza, debolezze, opportunità e minacce, fornendo una visione strategica completa.
    \item Un’altra parte cruciale è rappresentata dal \textit{business model}, che descrive come l’impresa intende generare valore e profitti, accompagnato dalla \textit{strategia di marketing e vendite}. 
    \item La sezione sulla \textit{roadmap} evidenzia le tappe fondamentali e le scadenze per il raggiungimento degli obiettivi.
    \item Non può mancare una descrizione dettagliata del \textit{team}, che illustra le competenze e le esperienze dei membri chiave dell’impresa, conferendo credibilità al progetto. 
    \item Infine, il \textit{piano economico-finanziario} fornisce proiezioni dettagliate di ricavi, costi e flussi di cassa, mentre la sezione sulla \textit{richiesta e utilizzo dei fondi} spiega come verranno investite le risorse eventualmente raccolte.
\end{enumerate}

\subsection*{Executive summary}
L’\textit{Executive Summary} è una delle sezioni più importanti di un business plan. Essendo la prima parte letta dagli investitori, rappresenta un vero e proprio biglietto da visita del progetto. Il suo scopo principale è catturare l’attenzione e invogliare il lettore a proseguire nella consultazione del documento. Per questo motivo, deve essere redatto in modo chiaro, sintetico ed efficace, racchiudendo gli elementi essenziali del piano imprenditoriale.

In un buon Executive Summary vengono presentati gli obiettivi principali dell’azienda, le strategie chiave, una sintesi delle previsioni finanziarie, i prodotti o servizi offerti e i fattori che rendono il progetto unico e degno di supporto. Ad esempio, se il business si basa su un’innovazione tecnologica o su un modello operativo particolarmente efficiente, questi aspetti devono emergere in modo immediato. Si consiglia di redigere questa sezione per ultima, dopo aver completato e analizzato ogni altra parte del piano, per garantire una visione d’insieme coerente e ben strutturata.

\subsection*{Mercato di riferimento}
La definizione del \textit{mercato di riferimento} è un passaggio fondamentale per comprendere a chi si rivolge il progetto imprenditoriale e quale sia il potenziale di sviluppo. Questa sezione deve fornire una descrizione chiara del pubblico target e del contesto competitivo, supportata da dati attendibili provenienti da fonti terze. Per esempio, l’analisi potrebbe includere statistiche sul numero di potenziali clienti, il volume d’affari del settore e i trend di crescita.

Un approccio comune prevede la suddivisione del mercato in tre categorie principali:
\begin{enumerate}
    \item \textit{TAM (Total Addressable Market)}, ovvero la domanda totale per un determinato prodotto o servizio.
    \item \textit{SAM (Served Addressable Market)}, il segmento di mercato effettivamente raggiungibile dall’azienda.
    \item \textit{SOM (Serviceable Obtainable Market)}, il mercato che realisticamente si prevede di conquistare.
\end{enumerate}

\subsection*{Problema}
Una delle basi per un business plan di successo è la descrizione chiara e dettagliata del \textit{problema} che si intende risolvere. In questa sezione, bisogna identificare le difficoltà affrontate dai potenziali clienti e spiegare perché le soluzioni esistenti non sono sufficienti. 

Per rendere questa analisi più incisiva, è utile includere dati economici che quantifichino il problema. 

\subsection*{Soluzione}
Dopo aver descritto il problema, la sezione successiva deve illustrare la \textit{soluzione} proposta, mettendone in evidenza i benefici concreti. Una buona descrizione della soluzione non deve essere eccessivamente tecnica, ma deve risultare chiara e comprensibile anche per chi non ha competenze specifiche nel settore.

Utilizzare casi pratici o scenari di utilizzo (use cases) aiuta a rendere la spiegazione più concreta e immediata.

\subsection*{Unique Selling Proposition (USP)}
La \textit{Unique Selling Proposition} è ciò che rende la proposta unica e convincente. Questa sezione del business plan deve spiegare chiaramente perché l’azienda o il prodotto si distinguono dalla concorrenza e in che modo migliorano la vita dei clienti. La USP è spesso il risultato di un mix di innovazione, efficienza operativa e capacità di soddisfare un bisogno specifico in modo migliore rispetto agli altri.

\subsection*{Competitor e vantaggi competitivi}
Un’analisi accurata della \textit{concorrenza} è essenziale per dimostrare la solidità della propria proposta di valore. Questa sezione deve identificare sia i competitor diretti, ovvero coloro che offrono prodotti o servizi simili, sia quelli indiretti, che soddisfano il medesimo bisogno attraverso soluzioni diverse.

Per distinguersi, è importante evidenziare i \textit{vantaggi competitivi}. 
Un’ottima pratica è raccogliere informazioni direttamente dai clienti o fornitori, o simulare l’esperienza utente richiedendo preventivi ai concorrenti. Questo approccio non solo arricchisce l’analisi, ma permette anche di comprendere meglio le aree su cui puntare per creare un vantaggio competitivo sostenibile.

\subsection*{SWOT analysis}
La \textit{SWOT Analysis} è uno strumento fondamentale per valutare la posizione competitiva di un’azienda e per pianificare strategie future. Questa analisi prende in considerazione quattro aree principali: punti di forza (Strengths), debolezze (Weaknesses), opportunità (Opportunities) e minacce (Threats).
\begin{itemize}
    \item I punti di forza di un’azienda possono includere competenze specifiche, tecnologie proprietarie o la presenza di risorse specializzate.
    \item Le debolezze, invece, rappresentano gli aspetti interni che limitano la competitività, come risorse finanziarie insufficienti o una proposta di valore non completamente chiara.
    \item Le opportunità riguardano i contesti esterni che l’azienda può sfruttare, come mercati emergenti o una domanda crescente per un determinato servizio. 
    \item Le minacce, infine, comprendono fattori esterni potenzialmente dannosi, come l’aumento della competizione o cambiamenti normativi sfavorevoli.
\end{itemize}

\subsection*{Business model, strategia di marketing e vendite}
Il \textit{business model} definisce come un’azienda genera valore e raggiunge i suoi clienti. Esistono diversi modelli applicabili, tra cui il SAAS (Software as a Service), il modello freemium, il marketplace transazionale o il noleggio. La scelta dipende dal tipo di prodotto o servizio e dal pubblico di riferimento, che può essere composto da aziende (B2B), consumatori finali (B2C) o una combinazione di entrambi (B2B2C).

Una \textit{strategia di marketing} ben definita è cruciale per il successo. Questa può includere pubblicità online, collaborazioni con influencer o eventi promozionali. Allo stesso modo, la \textit{strategia di vendita} varia in base al modello scelto.

\subsection*{Roadmap}
La \textit{roadmap} è la rappresentazione visiva di come l’azienda intende sviluppare il proprio prodotto e raggiungere i suoi obiettivi. Deve includere una timeline con le principali tappe previste, come il completamento di prototipi, il rilascio di versioni iniziali del prodotto o l’ottenimento di certificazioni cruciali.

\subsection*{Team}
Un \textit{team} competente e diversificato è spesso il fattore determinante per il successo di una startup. Ogni membro del team dovrebbe avere competenze complementari e un ruolo ben definito. 

Includere le foto dei membri del team e una breve descrizione delle loro esperienze può rendere questa sezione più personale e convincente.

\subsection*{Piano economico-finanziario}
Le \textit{proiezioni economico-finanziarie} a cinque anni sono essenziali per dimostrare la sostenibilità del progetto. Devono includere stime su ricavi, costi e margine operativo (EBITDA), preferibilmente presentate attraverso grafici che facilitano la comprensione.

I Key Performance Indicators (KPI) come il numero di utenti, i clienti acquisiti o i paesi raggiunti, sono strumenti utili per monitorare i progressi.

\subsection*{Richiesta e utilizzo dei fondi}
Per convincere gli investitori, è fondamentale presentare una \textit{richiesta} chiara dei fondi necessari e spiegare come verranno utilizzati nel tempo.
È altrettanto importante mostrare quali risultati si prevede di ottenere con questi fondi.

\section{Dossier}
Un \textit{dossier}\index{Business plan!dossier} ben strutturato è essenziale per presentare un progetto imprenditoriale agli investitori. La sua composizione deve seguire una sequenza logica e chiara, al fine di catturare l’interesse e dimostrare la serietà dell’iniziativa.
\begin{enumerate}
    \item \textit{Executive Summary:} Questa è la sezione iniziale, di una sola pagina o slide, che sintetizza i punti chiave del progetto.
    \item \textit{Teaser:} Una sezione introduttiva più approfondita, composta da tre pagine o slide, pensata per attirare l’attenzione e stimolare curiosità.
    \item \textit{Business Plan – Investor Deck:} Il cuore del dossier, suddiviso in 10-12 slide, che presenta in dettaglio i vari aspetti del progetto, dalla proposta di valore alla strategia economica.
    \item \textit{Relazione di accompagnamento:} Un documento narrativo, esteso tra le 15 e le 30 pagine, che approfondisce i contenuti del business plan.
    \item \textit{Proiezioni economico-finanziarie:} Fornite in un file Excel, queste proiezioni illustrano i ricavi, i costi e la sostenibilità del progetto nel medio-lungo termine.
    \item \textit{CV del team di progetto:} I curriculum dei componenti principali, evidenziando le esperienze e competenze rilevanti.
\end{enumerate}

\subsection{Guida Sequoia al Pitch Deck}
La \textit{guida Sequoia}\index{Business plan!guida sequoia} offre un approccio strutturato per creare un pitch deck efficace, suddiviso in vari elementi fondamentali.
\begin{itemize}
    \item \textit{Scopo della società.} Il pitch inizia con una dichiarazione chiara e concisa che descrive lo scopo dell’azienda. Questo richiede una sintesi incisiva che catturi l’essenza del progetto senza perdersi in dettagli tecnici o descrizioni prolisse.
    \item \textit{Il problema.} Descrivere il problema che i clienti affrontano è cruciale. Questo va fatto mostrando esempi concreti. È utile includere dati o storie per contestualizzare il problema.
    \item \textit{La soluzione.} La soluzione deve essere illustrata in modo semplice ma incisivo, con esempi pratici o use case.
    \item \textit{Perché ora?} Questa sezione esplora il contesto storico e le tendenze attuali che rendono possibile la soluzione proposta.
    \item \textit{Potenziale di mercato.} Definire chiaramente il pubblico di riferimento e il mercato è essenziale. La suddivisione in TAM, SAM e SOM aiuta a quantificare l’opportunità.
    \item \textit{Concorrenza e vantaggi competitivi.} Analizzare i concorrenti, diretti e indiretti, aiuta a dimostrare la consapevolezza del mercato.
    \item \textit{Prodotto.} Descrivere il prodotto e il suo sviluppo, dall’architettura alle funzionalità uniche, è fondamentale. Includere immagini o diagrammi del prodotto, come un’interfaccia intuitiva, può migliorare la comprensione.
    \item \textit{Modello di business.} Il modello di business deve chiarire come l’azienda genera ricavi e scala.
    \item \textit{Team.} La composizione del team è cruciale, specialmente nelle fasi iniziali. Ogni membro deve essere presentato con esperienze e competenze che lo rendono adatto a risolvere il problema.
    \item \textit{Finanziamenti e visione.} Le informazioni finanziarie, come profitti e perdite o flussi di cassa, rafforzano la credibilità del progetto. Infine, la visione a lungo termine deve ispirare, mostrando cosa sarà costruito nei prossimi cinque anni.
\end{itemize}

\subsection{Modello di Y Combinator}
La \textit{guida Y Combinator}\index{Business plan!Y combinator} segue un approccio simile, enfatizzando l’importanza della chiarezza e dell’impatto visivo.
\begin{itemize}
    \item \textit{Pagina del titolo.} La pagina iniziale deve presentare il nome dell’azienda e una descrizione sintetica della sua attività in una singola frase. Questa frase dovrebbe essere chiara, incisiva e capace di spiegare immediatamente il valore dell’impresa.
    \item \textit{Il problema.} Questa sezione è dedicata alla descrizione del problema che l’azienda intende risolvere. È importante spiegare chiaramente quale difficoltà i clienti affrontano e in che modo questo problema impatta la loro vita o le loro attività. Utilizzare dati concreti, aneddoti o esempi reali aiuta a rendere la questione tangibile.
    \item \textit{La soluzione.} Qui si spiega in modo chiaro e conciso come l’azienda risolve il problema. La descrizione deve concentrarsi sui benefici concreti e immediati per il cliente, evitando tecnicismi e dettagli eccessivi.
    \item \textit{La trazione.} Se l’azienda ha già ottenuto risultati significativi, questa è l’occasione per mostrarli. Può trattarsi di numeri relativi a utenti, ricavi, partnership, o altre metriche chiave. I grafici e le statistiche visive sono particolarmente efficaci in questa sezione.
    \item \textit{Altre metriche.} Se ci sono ulteriori dati rilevanti, come il tasso di conversione, la retention o metriche specifiche di settore, questa è la sezione in cui evidenziarli. Questi numeri aiutano a costruire una narrativa di successo e potenziale crescita.
    \item \textit{Proposta Unica di Valore (USP).} Questa sezione deve spiegare cosa rende l’azienda unica rispetto alla concorrenza. È il momento di evidenziare intuizioni strategiche, vantaggi tecnologici o altre caratteristiche distintive.
    \item \textit{Modello di business.} Il modello di business descrive come l’azienda genera ricavi e scala le operazioni. È utile dettagliare le principali fonti di reddito, il pricing e i segmenti di clientela.
    \item \textit{Il mercato e la crescita.} Questa sezione deve convincere gli investitori che il mercato è ampio e in espansione, e che l’azienda è ben posizionata per catturarne una quota significativa. È importante utilizzare dati di mercato, proiezioni di crescita e spiegare come l’azienda intende espandersi.
    \item \textit{Il team.} Gli investitori attribuiscono grande importanza al team, specialmente nelle fasi iniziali. Questa sezione deve presentare i fondatori e i membri chiave, evidenziando le loro competenze ed esperienze pertinenti.
    \item \textit{La richiesta di fondi.} Infine, è importante indicare chiaramente quanto capitale l’azienda sta cercando, come verrà utilizzato e quali risultati saranno raggiunti grazie all’investimento.
\end{itemize}


\chapter{Piano finanziario}
Un'idea innovativa, per quanto brillante, non può decollare senza un piano finanziario solido che ne garantisca la sostenibilità economica. Questo strumento è cruciale per rispondere a domande fondamentali che ogni imprenditore dovrebbe porsi: quanto denaro serve per sviluppare e lanciare il progetto? Quali saranno i ricavi e i costi nel tempo? E soprattutto, quando l’azienda inizierà a generare profitti?

\begin{redbar}
Il \textbf{piano finanziario}\index{Business plan!piano finanziario} fornisce una mappa dettagliata per gestire l’aspetto economico di un’attività. Per esempio, aiuta a calcolare l’investimento iniziale e il capitale necessario per coprire i costi operativi nei primi mesi o anni. Inoltre, offre previsioni sui guadagni e sulle spese, evidenziando se e quando il progetto diventerà sostenibile. Un elemento cruciale è il break-even point, ovvero il momento in cui i ricavi superano i costi, dimostrando la capacità del progetto di generare utili.
\end{redbar}

\section{Struttura del piano finanziario}
Il piano finanziario si articola in tre componenti principali: il budget iniziale, il break-even point e le proiezioni finanziarie.

\subsection{Budget iniziale}
La prima fase, chiamata \textit{budget iniziale}\index{Business plan!budget iniziale}, consiste nell’identificazione dei costi necessari per avviare l’attività. Questi includono:
\begin{itemize}
    \item \textit{Costi tecnologici}, come server, licenze software e hardware.
    \item \textit{Costi operativi}, quali affitto, utenze e spese amministrative.
    \item \textit{Costi di marketing}, come campagne pubblicitarie e ottimizzazione SEO.
    \item \textit{Costi di sviluppo}, che comprendono stipendi, prototipi e test.
\end{itemize}
Un esempio pratico potrebbe essere il lancio di una piattaforma di tutoraggio online. I costi iniziali comprenderebbero 3.000 euro per un anno di server cloud (ad esempio, Microsoft Azure), 40.000 euro per il lavoro di due sviluppatori per quattro mesi e 5.000 euro per una campagna di marketing iniziale. Il capitale totale necessario ammonta quindi a 3.000 euro + 40.000 euro + 5.000 euro = 48.000 euro.

\subsection{Proiezioni finanziarie}
Le \textit{proiezioni finanziarie}\index{Business plan!proiezioni finanziarie} offrono un quadro completo di come l’attività genererà ricavi e coprirà i costi nel tempo. Si sviluppano attraverso tre strumenti principali: il conto economico, il flusso di cassa e lo stato patrimoniale.
\begin{enumerate}
    \item Il \textit{conto economico} analizza ricavi, costi fissi, costi variabili e l’utile o la perdita d’esercizio.
    \begin{equation*}
        \textit{Utile Netto} = \textit{Ricavi} - (\textit{Costi Fissi} + \textit{Costi Variabili})
    \end{equation*}
    Ad esempio, per un’app mobile con un abbonamento mensile di 10 euro e 500 utenti, i ricavi mensili sono di 5.000 euro. Se i costi fissi ammontano a 3.000 euro (per server e personale) e i costi variabili a 1.000 euro (banda e assistenza clienti), l’utile netto mensile sarà di 5.000 euro - (3.000 euro + 1.000 euro) = 1.000 euro.

    \item Il \textit{flusso di cassa} misura la liquidità dell’azienda, mostrando i movimenti effettivi di denaro in entrata e in uscita. A differenza dell’utile netto, che include valori contabili come ammortamenti, il flusso di cassa riflette esclusivamente il denaro reale. 
    \begin{equation*}
        \textit{Cash Flow} = \textit{Entrate} - \textit{Uscite}
    \end{equation*}
    Per esempio, un’attività con entrate di 12.000 euro (da abbonamenti e servizi) e uscite di 9.000 euro (per server e stipendi) genera un flusso di cassa positivo di 12.000 euro - 9.000 euro = 3.000 euro.

    \item Lo \textit{stato patrimoniale} fornisce una fotografia delle risorse e dei debiti di un’azienda. Comprende:
    \begin{itemize}
        \item L’attivo, che include risorse come server, software sviluppato e liquidità disponibile.
        \item Il passivo, che rappresenta debiti e finanziamenti ricevuti.
        \item Il patrimonio netto, ovvero la differenza tra attivo e passivo.
    \end{itemize}
    \begin{equation*}
        \textit{Patrimonio Netto} = \textit{Attivo} - \textit{Passivo}
    \end{equation*}
    Un esempio pratico mostra un attivo totale di 100.000 euro (30.000 euro per server, 50.000 euro per software e 20.000 euro di liquidità) e un passivo di 40.000 euro (30.000 euro di prestiti e 10.000 euro di debiti verso fornitori). Il patrimonio netto è quindi di 100.000 euro - 40.000 euro = 60.000 euro.
\end{enumerate}

\subsection{Break-even point}
Il \textit{break-even point}\index{Business plan!break-even point} rappresenta il momento in cui i ricavi dell’azienda coprono interamente i costi, segnando l’inizio della generazione di utili. Si calcola dividendo i costi fissi per il margine di contribuzione (cioè il prezzo unitario meno il costo variabile per unità). 
\begin{equation*}
    \textit{Break-even Point} = \frac{\textit{Costi Fissi}}{\textit{Prezzo unitario} – \textit{Costo Variabile per unità}}
\end{equation*}
Ad esempio, per un software venduto a 100 euro per licenza, con costi fissi annuali di 20.000 euro e un costo variabile per licenza di 20 euro, il break-even point si raggiunge con la vendita di 20.000 euro/(100 euro - 20 euro) = 250 licenze.

\section{Costruzione di un piano finanziario}
Un piano finanziario solido richiede l’identificazione accurata di tutti i costi e una proiezione realistica dei ricavi. Questi elementi devono essere analizzati con attenzione per garantire una gestione efficace delle risorse e per dimostrare la sostenibilità economica del progetto.

\subsection{Identificare i costi}
I costi di un progetto si suddividono in due categorie principali: costi fissi e costi variabili.
\begin{itemize}
    \item I \textit{costi fissi} rappresentano le spese che rimangono costanti indipendentemente dal volume di attività. Tra questi, si includono l’affitto e le spese generali, come bollette per luce e acqua, che sono essenziali per gestire uno spazio fisico o un ufficio. Anche gli stipendi del personale, come sviluppatori, project manager e amministratori, rientrano in questa categoria. Un altro elemento fondamentale sono i costi per licenze software, come strumenti di sviluppo e piattaforme di gestione, oltre a spese legali e assicurative necessarie per proteggere l’attività.

    Ad esempio, avviando una piattaforma SaaS (Software as a Service) per la gestione delle risorse aziendali, ipotizzando di acquisire 500 utenti: i costi fissi possono essere stimati in 1.500 euro per l’affitto dei locali, 12.000 euro per il personale e 2.000 euro per le licenze software, per un totale di 15.500 euro.

    \item I \textit{costi variabili} dipendono dal volume di attività e aumentano proporzionalmente con la crescita del progetto. Per una piattaforma SaaS, i principali costi variabili includono le spese per banda e hosting, che crescono con il numero di utenti, e i costi di supporto clienti. Anche le spese di marketing, come campagne pubblicitarie sono variabili, poiché dipendono dalle dimensioni e dalla durata delle campagne.

    Nel caso della piattaforma SaaS con 500 utenti, ogni utente genera un costo di 2 euro per il cloud e 1 euro per il supporto clienti, per un totale di 3 euro per utente. Con 500 utenti, i costi variabili ammontano a 1.500 euro, portando i costi totali (fissi + variabili) a 15.500 euro + (500 euro $\times$ 3 euro) = 17.000 euro.
\end{itemize}

\subsection{Stimare i ricavi}
La stima dei ricavi è altrettanto cruciale per valutare la sostenibilità di un progetto. I ricavi possono provenire da diverse fonti, tra cui abbonamenti, pubblicità e vendita di prodotti o servizi aggiuntivi.

Per un progetto SaaS con 500 utenti che pagano 30 euro al mese, i ricavi mensili generati dagli abbonamenti sono pari a 15.000 euro. A questi si possono aggiungere entrate dalla pubblicità, come la vendita di spazi pubblicitari o il pay-per-click, che potrebbe generare 500 euro. Infine, la vendita di corsi online o consulenze premium potrebbe portare ulteriori 3.000 euro, per un totale di 15.000 euro + 500 euro + 3.000 euro = 18.500 euro di ricavi mensili.

\subsection{Costruzione delle proiezioni finanziarie}
Un piano finanziario efficace non si limita a una singola previsione statica ma include proiezioni dinamiche su più scenari. Le proiezioni si estendono tipicamente su 3-5 anni e includono tre scenari principali:
\begin{itemize}
    \item Scenario realistico, che considera una crescita moderata e costante.
    \item Scenario ottimistico, basato su una crescita accelerata grazie a un forte marketing o a fattori virali.
    \item Scenario pessimistico, che analizza le conseguenze di vendite inferiori alle aspettative, dovute a concorrenza o difficoltà di mercato.
\end{itemize}
Per supportare queste proiezioni, è essenziale costruire un conto economico e un flusso di cassa operativo per i prossimi cinque anni, includendo tutte le variabili possibili.

\subsection{EBIT ed EBITDA}
Due indicatori chiave per comprendere la redditività operativa di un progetto sono l’EBIT e l’EBITDA.
\begin{itemize}
    \item L’\textit{EBIT (Earnings Before Interest and Taxes)} misura l’utile operativo dell’azienda prima di considerare interessi e imposte. Questo indicatore consente di valutare se l’impresa è profittevole dal punto di vista operativo, escludendo gli effetti delle scelte fiscali o della struttura del debito.
    \begin{equation*}
        \textit{EBIT} = \textit{Ricavi} – \textit{Costi operativi (esclusi interessi e tasse)}
    \end{equation*}
    
    \item L’\textit{EBITDA (Earnings Before Interest, Taxes, Depreciation, and Amortization)} esclude anche ammortamenti e svalutazioni, offrendo una visione più chiara della capacità dell’impresa di generare cash flow operativo. Se l’EBITDA è negativo, l’azienda sta perdendo denaro anche prima di considerare gli investimenti.
    \begin{equation*}
        \textit{EBITDA} = \textit{EBIT} + \textit{Ammortamenti e Svalutazioni}
    \end{equation*}
\end{itemize}

\subsection{KPI}
I \textit{KPI (Key Performance Indicators)} sono strumenti fondamentali per misurare l'efficacia di un'azienda, un progetto o una strategia. Grazie ai KPI, è possibile monitorare e valutare in modo concreto il livello di successo rispetto agli obiettivi prefissati. Questi indicatori offrono una visione chiara delle performance aziendali, rendendoli essenziali per il processo decisionale e per la pianificazione strategica.

In generale, i KPI si dividono in due categorie principali:
\begin{itemize}
    \item \textit{KPI Finanziari:} Misurano il valore economico e la redditività dell’impresa.
    \item \textit{KPI Non Finanziari:} Analizzano aspetti legati alla soddisfazione del cliente, all’efficienza operativa e ad altri fattori qualitativi.
\end{itemize}
I KPI aiutano le aziende in diversi modi:
\begin{itemize}
    \item Monitorare la performance: Verificano se gli obiettivi aziendali sono stati raggiunti.
    \item Prendere decisioni informate: Forniscono dati concreti e oggettivi per supportare scelte strategiche.
    \item Pianificare il futuro: Evidenziano aree di miglioramento e offrono spunti per ottimizzare le operazioni aziendali.
\end{itemize}
Poiché ogni KPI analizza un aspetto specifico, è necessario considerare un insieme diversificato di indicatori per ottenere una visione completa delle performance aziendali.

Tra i KPI finanziari più utilizzati troviamo:
\begin{itemize}
    \item \textit{Gross Margin (Margine Lordo).} Il Gross Margin misura la redditività di un prodotto o servizio prima di considerare i costi operativi. Si calcola tenendo conto solo dei costi diretti, offrendo un’indicazione della capacità dell’azienda di generare profitti dalle vendite.
    \begin{equation*}
    \textit{Gross Margin (\%)} = \frac{\textit{Ricavi totali} - \textit{Costi diretti}}{\textit{Ricavi totali}}\times 100
    \end{equation*}
    Un margine lordo elevato indica che l’azienda trattiene una percentuale significativa dei ricavi come profitto prima di considerare i costi operativi.
    
    \item \textit{Operating Margin (Margine Operativo).} L’Operating Margin misura la redditività dopo aver sottratto tutti i costi operativi (sia fissi che variabili) dai ricavi totali. Questo indicatore mostra l’efficienza complessiva dell’azienda nel generare profitti operativi.
    \begin{equation*}
    \textit{Operating Margin(\%)} = \frac{\textit{Ricavi totali} - \textit{Costi operativi totali}}{\textit{Ricavi totali}}\times 100        
    \end{equation*}
    Un margine operativo positivo indica che l’azienda sta generando un profitto netto dalle operazioni.
    
    \item \textit{ROI (Return on Investment).} Il ROI valuta il rendimento di un investimento rispetto al capitale impiegato, offrendo una misura dell’efficienza e della profittabilità delle risorse investite.
    \begin{equation*}
    \textit{ROI (\%)} = \frac{\textit{Guadagno netto dell’investimento} - \textit{Costo dell’investimento}}{\textit{Costo dell’investimento}}\times 100
    \end{equation*}
    Questo indicatore è particolarmente utile per confrontare diverse opportunità di investimento e per valutare il successo di un progetto specifico.
\end{itemize}
Accanto ai KPI finanziari, gli indicatori non finanziari completano il quadro fornendo informazioni qualitative sulla performance dell’azienda. I principali KPI non finanziari includono:
\begin{itemize}
    \item \textit{ARPU (Average Revenue Per User).} L’ARPU misura il ricavo medio generato da ogni cliente o utente, aiutando a comprendere il valore economico di ciascun cliente per l’azienda. Si calcola dividendo le entrate totali per il numero di clienti totali:
    \begin{equation*}
        \textit{ARPU} = \frac{\textit{Entrate totali}}{\textit{Numero di clienti totali}}
    \end{equation*}
    Un ARPU elevato suggerisce che i clienti generano un valore maggiore per l’azienda, indicando una qualità superiore del prodotto o servizio.
    \item \textit{Churn Rate.} Il Churn Rate misura la percentuale di clienti che abbandonano un servizio o cancellano un abbonamento in un determinato periodo. È un indicatore cruciale per valutare la soddisfazione del cliente e la fidelizzazione.
    \begin{equation*}
    \textit{Churn Rate (\%)} = \frac{\textit{Clienti persi nel mese}}{\textit{Clienti totali all’inizio del mese}}\times 100
    \end{equation*}
    Un Churn Rate elevato può indicare problemi con il prodotto, il servizio o l’esperienza del cliente, richiedendo interventi mirati per migliorare questi aspetti.
    \item \textit{MRR (Monthly Recurring Revenue).} Il MRR rappresenta le entrate ricorrenti mensili generate da abbonamenti o contratti a lungo termine. È particolarmente rilevante per i modelli di business basati su entrate continuative, come le piattaforme SaaS.
    \begin{equation*}
    \textit{MRR}=\textit{Numero di clienti paganti}\times\textit{Prezzo medio mensile dell’abbonamento}   
    \end{equation*}
    Un MRR in crescita è un segnale positivo per la stabilità finanziaria e il successo di lungo termine dell’azienda.
\end{itemize}

\section{Case study}
Uno degli esempi più emblematici di business plan ben strutturato e di successo in Italia è quello di Satispay, una startup fintech che ha rivoluzionato il mondo dei pagamenti digitali. Fondata nel 2013, Satispay ha sviluppato un'applicazione per smartphone che consente transazioni rapide, sicure e senza la necessità di carte di credito, offrendo un'alternativa semplice e conveniente sia per privati sia per piccole e medie imprese.
\begin{enumerate}
    \item \textit{Analisi di mercato.} Satispay nasce dall’identificazione di un’esigenza specifica: ridurre le elevate commissioni sui pagamenti digitali che gravavano sui piccoli commercianti e superare le complessità dell'infrastruttura bancaria tradizionale. Questo insight ha portato i fondatori a progettare una soluzione low-cost che fosse facilmente accessibile da un mercato in rapida crescita. L'obiettivo era chiaro: democratizzare i pagamenti digitali rendendoli semplici e convenienti per tutti.
    \item \textit{Modello di business.} Il modello di business di Satispay si basa su un principio fondamentale: la gratuità per i privati e commissioni fisse, ma ridotte, per gli esercenti. Questo approccio, oltre a fidelizzare gli utenti privati, ha reso l'app particolarmente attraente per i piccoli commercianti, spesso penalizzati dai costi elevati dei sistemi di pagamento tradizionali. Con questa strategia, Satispay ha potuto scalare rapidamente, puntando a una base utenti ampia e attrattiva anche per i grandi marchi.
    \item \textit{Innovazione tecnologica.} Uno degli elementi distintivi di Satispay è l’utilizzo del codice IBAN per bypassare l’infrastruttura delle carte di credito. Questo approccio ha permesso di offrire un sistema di pagamento tecnologicamente semplice ma sicuro, investendo in funzionalità che hanno costruito fiducia nei consumatori. La sicurezza e l’usabilità sono stati due pilastri fondamentali per il successo dell’app.
    \item \textit{Strategia di marketing.} Fin dal principio, Satispay ha adottato una strategia di marketing mirata, collaborando con esercenti locali per promuovere l'app come una soluzione conveniente. Campagne di cashback e programmi di referral hanno accelerato l’acquisizione di utenti, sfruttando il passaparola come motore di crescita. Questo approccio ha creato un effetto rete che ha rapidamente aumentato la base utenti e il numero di esercenti aderenti.
    \item \textit{Sostenibilità finanziaria.} Un aspetto cruciale del successo di Satispay è stata la capacità di attrarre finanziamenti significativi. Dopo un investimento iniziale di circa 400.000 euro raccolti dai fondatori e da alcuni business angel, la startup ha condotto numerosi round di raccolta fondi. Tra i più rilevanti, nel 2015 ha ottenuto 5 milioni di euro da investitori strategici come Iccrea Banca, seguito da un round da 18,3 milioni di euro nel 2017. Nel 2022, Satispay ha raccolto ben 320 milioni di euro, raggiungendo una valutazione di oltre 1 miliardo di euro, diventando così uno dei pochi "unicorni" italiani.
\end{enumerate}
Il successo della startup è testimoniato dai numeri. Nel periodo iniziale, tra il 2013 e il 2015, Satispay si era posta obiettivi ambiziosi, come raggiungere 100.000 utenti attivi e coinvolgere 5.000 esercenti. Con un volume di transazioni annuale superiore a 1 miliardo di euro nel 2022 e una base di 3 milioni di utenti attivi, l’azienda ha superato ampiamente queste previsioni. Il 2022 ha segnato anche il raggiungimento di 200.000 esercenti aderenti in Italia e in Europa, confermando la validità del modello di business.

Guardando al futuro, Satispay punta a espandersi ulteriormente in Europa, con particolare attenzione a mercati come Francia e Germania. L’obiettivo dichiarato è di raggiungere 5 milioni di utenti entro il 2024, aumentando il volume delle transazioni del 50\% annuo. Nel frattempo, l’azienda ha dichiarato di essere vicina alla sostenibilità operativa, grazie alle entrate generate dalle commissioni fisse e dagli accordi con grandi marchi.

Il caso di Satispay dimostra come un business plan ben strutturato, unito a un’esecuzione efficace, possa trasformare una startup in un successo internazionale. Tra le lezioni chiave emergono:
\begin{itemize}
    \item \textit{Chiarezza degli obiettivi:} L’azienda ha saputo identificare una necessità specifica nel mercato e proporre una soluzione scalabile ed efficace.
    \item \textit{Eccellenza nell'implementazione tecnologica:} L'utilizzo del codice IBAN ha reso il sistema innovativo e facilmente adottabile.
    \item \textit{Efficacia nella raccolta fondi:} Attraverso round progressivi, Satispay ha attratto investitori grazie a un piano solido e a una roadmap di crescita convincente.
\end{itemize}
